\documentclass[conference]{IEEEtran}

\usepackage{graphicx}
\usepackage{multirow}
\usepackage{array}
\usepackage{caption}
\usepackage{subcaption}
\usepackage{xcolor}
\usepackage{colortbl}

\usepackage{amsmath, amssymb, amsthm}
\usepackage{bm}
\usepackage{mathtools}
\usepackage{algorithm}
\usepackage{algpseudocode}
\usepackage{booktabs}

\usepackage{graphicx,color}
\newcommand{\blue}[1]{{\color{blue}{#1}}}
\newcommand{\red}[1]{{\color{red}{#1}}}
\newcommand{\mbf}[1]{\mathbf{#1}}
\newcommand{\abs}[1]{\left|#1\right|}
\newcommand{\norm}[1]{\left\|{#1}\right\|}

\newcommand{\EE}{\mathbb{E}}
\newcommand{\RR}{\mathbb{R}}
\DeclareMathOperator*{\argmin}{arg\,min}
\DeclareMathOperator*{\argmax}{arg\,max}
\newtheorem{remark}{Remark}

\def\BibTeX{{\rm B\kern-.05em{\sc i\kern-.025em b}\kern-.08em
    T\kern-.1667em\lower.7ex\hbox{E}\kern-.125emX}}

\begin{document}

\title{Post-Decision-State Proximal Policy Optimization for Constrained UAV 360\textdegree{} Video Streaming: A PDS-GAE Extension and Empirical Characterization}

\author{
\IEEEauthorblockN{Jad Al Lakkis}
\IEEEauthorblockA{% TODO: affiliation, advisor(s) as co-authors, dept./university
\textit{[affiliation]} \\
[email]}
}

\maketitle

% =====================================================================
\begin{abstract}
We study tile-selective enhancement-layer streaming of 360\textdegree{} video from a UAV under
joint stall, transmit-power, and tile-budget constraints, formulated as a constrained
Markov decision process (CMDP) and solved with Lagrangian-relaxed proximal policy
optimization (PPO). Building on this baseline, we introduce a post-decision-state (PDS)
critic that separates the reward and next state into a component known immediately
after the tile/power decision and a component that remains random until the next
viewport and channel realization are observed. An initial one-step PDS advantage,
adapted from prior soft-actor-critic-based PDS work, is shown empirically to truncate
credit assignment and fail to reach reliable feasibility; we derive and evaluate a
generalization, PDS-GAE, that chains the PDS residual through the same backward
generalized-advantage-estimation recursion used by the baseline. Across two constraint
budgets, multiple seeds, and matched-rollout comparisons, PDS-GAE reaches moderate
constraint satisfaction $2$--$3.5\times$ faster than the PPO baseline in every seed
tested, an effect we trace to a measurably lower-variance policy-gradient residual
relative to a same-rollout ordinary-TD counterfactual. This early advantage does not,
however, translate into uniformly higher final reliability: PDS-GAE exhibits a
persistent, non-resolving feasibility oscillation, independently reproduced in both
budgets over $4{,}000+$ training rollouts, which we characterize and partially
diagnose. We further propose Virtual Experience (VE), a mechanism that reuses one
observed random outcome across multiple hypothetical actions to generate additional
training targets for the PDS critic at zero extra environment cost; we give its
derivation, algorithm, and applicability argument, but do not yet report empirical
results for it, leaving its evaluation to future work.
\end{abstract}

\begin{IEEEkeywords}
constrained reinforcement learning, proximal policy optimization, post-decision state,
360-degree video streaming, UAV communications, Lagrangian relaxation
\end{IEEEkeywords}

% =====================================================================
\section{Introduction}

Unmanned aerial vehicles (UAVs) equipped with $360^{\circ}$ cameras are an attractive
platform for immersive live streaming, but the wireless downlink to a head-mounted
display (HMD) is simultaneously bandwidth-, power-, and buffer-limited. A standard
mitigation, used in~\cite{ref_mmsp25} % TODO: verify exact MMSP'25 citation key/venue
and adopted as the base system here, is to encode the panorama as a low-resolution
base layer plus a tile-partitioned enhancement layer, and to transmit only the
enhancement-layer tiles that overlap the user's \emph{predicted} viewport, plus any
additional tiles an agent chooses to send speculatively. This is naturally posed as a
constrained Markov decision process (CMDP): maximize expected discounted viewport
quality subject to per-step stall-time, transmit-power, and tile-count budgets.
Lagrangian relaxation with dual-ascent multiplier updates converts this into an
unconstrained reward that a standard policy-gradient method, here PPO, can optimize
directly, as established in our base system~\cite{ref_mmsp25}.

Every quantity needed to compute the stall time, transmitted data volume, and
buffer evolution is a \emph{deterministic} function of the state and the chosen
action; only the realized viewport (used to score quality after the fact) and the
next wireless channel gain remain unknown at decision time. This known/random
split is exactly the \emph{post-decision state} (PDS) construction from approximate
dynamic programming~\cite{powell_adp}, previously combined with soft-actor-critic in
a similar streaming setting~\cite{ref_ehs}. % TODO: verify exact EHS citation
Intuitively, a critic that conditions on the deterministic post-decision quantities
before bootstrapping should be a Rao-Blackwellized -- and hence lower-variance --
estimate of the value function than one that bootstraps through the full, still-random
next state. This paper adapts that PDS mechanism to PPO's on-policy, clipped-surrogate
setting, distinct from a concurrent, independently-derived PPO+PDS method with different
advantage equations~\cite{felizardo2025pds}.

\textbf{Contributions.}
\begin{enumerate}
\item We derive \emph{PDS-GAE}, a generalization of the original one-step PDS advantage
that chains the PDS residual through the same backward $\mathrm{GAE}(\lambda)$
recursion~\cite{schulman2016gae} used by the PPO baseline. We show via an ablation
that the one-step estimator, not the second-critic architecture itself, was the
cause of an earlier PDS variant's failure to reach reliable feasibility.
\item We evaluate PDS-GAE against the PPO baseline across two constraint budgets,
multiple seeds, and matched training-rollout counts, using the same held-out
evaluation protocol for every run.
\item We give direct empirical evidence for PDS-GAE's core theoretical claim -- a
lower-variance advantage residual -- computed as a zero-extra-cost counterfactual on
the same rollouts, and connect it to an observed $2$--$3.5\times$ speedup in reaching
moderate constraint satisfaction.
\item We identify, and partially diagnose, a persistent feasibility oscillation in
PDS-GAE that is not present in the PPO baseline and does not resolve with substantially
more training, reporting the candidate mechanisms we tested and ruled out.
\item We propose \emph{Virtual Experience} (VE), which reuses one realized,
action-independent random outcome across several hypothetical actions to generate
additional PDS-critic training targets at no extra environment-interaction cost. We
give its algorithm and applicability argument but explicitly do not yet report results
for it.
\end{enumerate}

\emph{Notation.} Scalars are non-bold ($a$), vectors bold-lowercase ($\mbf{a}$).
$\EE[\cdot]$ is the expectation operator, $\gamma$ the discount factor, and
$[\cdot]_+ \triangleq \max(0,\cdot)$.

% =====================================================================
\section{Related Work}

\textbf{Base streaming system.} Our CMDP formulation, tile/action space, and
Lagrangian-PPO baseline follow~\cite{ref_mmsp25} directly; we extend it with a
post-decision-state critic and evaluate the extension under the same reward and
constraint structure.

\textbf{Post-decision states.} PDS variables are a long-standing approximate dynamic
programming tool for separating deterministic and random transition dynamics to
reduce value-estimation variance~\cite{powell_adp}. \cite{ref_ehs} combines PDS with
soft actor-critic in a related streaming setting; the one-step advantage we initially
adapt from that formulation is Eq.~(28)--(29) below. An independent, concurrent PPO+PDS
method~\cite{felizardo2025pds} uses different advantage equations and is not the
source of the PDS-GAE trace we derive in Sec.~\ref{sec:method}.

\textbf{Advantage estimation.} PDS-GAE applies the backward, $\lambda$-weighted
residual-chaining recursion of generalized advantage estimation~\cite{schulman2016gae}
to the PDS residual in place of the ordinary TD residual; the resulting advantage is
used inside PPO's clipped surrogate objective~\cite{schulman2017ppo}.

\textbf{Hindsight/counterfactual value estimation.} Virtual Experience's core idea --
reusing an outcome that is conditionally independent of the action, given the state,
to obtain lower-variance value information about actions not actually taken -- is a
member of a broader family that includes hindsight credit assignment~\cite{hca_harutyunyan}
% TODO: verify exact HCA citation/arXiv id before submission
and counterfactual credit assignment~\cite{mesnard2021cca}. Neither paper's mechanism
matches VE's concrete instantiation (recomputing deterministic physics for several
hypothetical actions against one reused random draw to generate extra post-decision
critic targets); we treat VE's novelty at the level of \emph{that specific mechanism}
as plausible but not yet verified against the full literature.

% =====================================================================
\section{System Model}
\label{sec:system}

We adopt the CMDP formulation of~\cite{ref_mmsp25}. The panorama is partitioned into
tiles $\mathcal{N}=\{1,\dots,N\}$. The predicted-viewport tiles are always sent; the
agent additionally selects any subset of the remaining tiles. The enhanced region and
action are
\begin{align}
\mathcal{E}^t(\mbf{x}^t) &= \{i \in \mathcal{N} : x_i^t = 1\}, \\
\alpha^t &= (\mbf{x}^t, P^t), \qquad 0 \le P^t \le P_{\max}.
\end{align}
The CMDP maximizes discounted viewport quality subject to stall-time, power, and
tile-count budgets:
\begin{align}
\max_\pi \; & \EE_\pi\Big[\textstyle\sum_{t} \gamma^t \bar Q^t_{\text{viewport}}\Big] \\
\text{s.t. } & \EE_\pi\Big[\textstyle\sum_t \gamma^t D^t / T_0\Big] \le \bar D, \\
& \EE_\pi\Big[\textstyle\sum_t \gamma^t P^t / P_{\max}\Big] \le \bar P, \\
& \EE_\pi\Big[\textstyle\sum_t \gamma^t \textstyle\sum_i x_i^t / N\Big] \le \bar B,
\end{align}
with $\gamma=0.99$. The penalty-based reward is
\begin{equation}
r^t = \beta_q \bar Q^t_{\text{viewport}} - \beta_s \tfrac{D^t}{T_0} - \beta_p \tfrac{P^t}{P_{\max}} - \beta_b \tfrac{\sum_i x_i^t}{N}.
\end{equation}
The pre-decision state is $\omega^t = (Z^t, \Delta_\theta^{t-1}, \Delta_\phi^{t-1}, h^t)$
(buffer level, previous viewport-prediction errors, current channel gain), and coverage
quality is $\bar Q^t_{\text{viewport}} = C^t_{\text{cov}} = |V^t_{\text{actu}} \cap
\mathcal{E}^t(\mbf{x}^t)| / |V^t_{\text{actu}}|$, with $V^t_{\text{actu}}$ the actual
viewport. Once the action is chosen, the enhanced region, transmitted data volume
$A^t$, transmission rate $R^t$ (Shannon capacity over the UAV--BS channel), stall time
$D^t$, and the deterministic buffer update $\tilde Z^{t+1}$ are all known; only the
actual viewport and the next channel gain $h^{t+1}$ remain random.

\textbf{Lagrangian relaxation.} Following~\cite{ref_mmsp25}, fixed weights
$\beta_s,\beta_p,\beta_b$ are replaced by adaptive multipliers
$\mu_D,\mu_P,\mu_B \ge 0$, updated once per rollout by projected dual ascent,
$\mu_j \leftarrow [\mu_j + \eta_\mu (\bar J_j^{\text{rollout}} - \bar b_j)]_+$,
$j\in\{D,P,B\}$, where $\bar J_j^{\text{rollout}}$ is the sample-average discounted
constraint return over the rollout's episodes.

% =====================================================================
\section{Method: PDS-GAE and Virtual Experience}
\label{sec:method}

\subsection{Post-decision state and the known/random reward split}

The post-decision state retains everything already determined by the action:
\begin{equation}
\tilde\omega^t = \big(\tilde Z^{t+1}, \Delta_\theta^{t-1}, \Delta_\phi^{t-1}, h^t, \mathcal{E}^t(\mbf{x}^t)\big).
\end{equation}
The reward splits accordingly into a known and a random part,
\begin{align}
r^t_{\text{known}} &= -\mu_D c_D^t - \mu_P c_P^t - \mu_B c_B^t, \label{eq:rknown}\\
r^t_{\text{random}} &= \beta_q C_{\text{cov}}^t, \label{eq:rrandom}
\end{align}
with $c_D^t = D^t/T_0$, $c_P^t = P^t/P_{\max}$, $c_B^t = \sum_i x_i^t / N$ the same three
constraint costs. A second critic $\tilde V_\psi(\tilde\omega^t)$, architecturally
independent of the ordinary critic $V_\phi(\omega^t)$, is trained toward the one-step
target
\begin{equation}
\tilde y^t = r^t_{\text{random}} + \gamma(1-d^t)\, V_{\text{old}}(\omega^{t+1}), \label{eq:ytilde}
\end{equation}
where $d^t{=}1$ at a true episode terminal and $V_{\text{old}}$ denotes the
pre-decision critic \emph{as evaluated while collecting the rollout} (frozen for the
policy-gradient step).

\subsection{The original one-step PDS advantage, and why it fails here}

Substituting $Q^\pi(\omega^t,\alpha^t) = r^t_{\text{known}} + \tilde V^\pi(\tilde\omega^t)$
into $A^\pi = Q^\pi - V^\pi$ gives the PDS residual adapted from~\cite{ref_ehs}:
\begin{equation}
\delta_t^{\text{pre}} = r_{\text{known}}^t + \tilde V(\tilde\omega^t) - V_{\text{old}}(\omega^t). \label{eq:delta_pre}
\end{equation}
Used directly as the actor's advantage (Run C in our internal ablation naming), this
one-step estimator truncates credit assignment to a single step. In a
constrained-Lagrangian environment where a late-episode constraint cost must propagate
back over a $36$-step episode, an ablation isolating the estimator from the critic
architecture (an ordinary single critic given the same one-step bootstrap, Run B)
showed the same qualitative failure as the PDS critic itself (Run C): both plateau with
two of three multipliers decaying to zero -- i.e., \emph{abandoning} rather than
satisfying those constraints -- while the GAE($0.95$) baseline (Run A) keeps all three
multipliers live. This isolates the one-step horizon, not the second critic, as the
cause, motivating the following generalization.

\subsection{PDS-GAE}

We chain $\delta_t^{\text{pre}}$ through the same backward,
$\lambda$-weighted recursion GAE~\cite{schulman2016gae} applies to the ordinary TD
residual:
\begin{equation}
\hat A_t^{\text{PDS-GAE}} = \delta_t^{\text{pre}} + \gamma\lambda_{\text{GAE}}(1-d_t)\,\hat A_{t+1}^{\text{PDS-GAE}}, \label{eq:pdsgae}
\end{equation}
with $\gamma{=}0.99$, $\lambda_{\text{GAE}}{=}0.95$ matching the baseline exactly, and
the recursion reset to zero at every true episode terminal. The pre-decision critic is
fitted to $\hat R_t^{\text{PDS-GAE}} = \hat A_t^{\text{PDS-GAE}} + V_{\text{old}}(\omega^t)$;
the PDS critic's own target, Eq.~\eqref{eq:ytilde}, is unchanged. At $\lambda_{\text{GAE}}{=}0$
this collapses exactly to the one-step advantage, Eq.~\eqref{eq:delta_pre}.

\begin{remark}
Under \emph{exact} on-policy critics, $\delta_t^{\text{pre}}$ already equals the true
advantage $A^\pi(\omega^t,\alpha^t)$, so later terms in Eq.~\eqref{eq:pdsgae} have
conditional mean zero given the present action -- the PDS-GAE trace does not inherit
ordinary GAE's telescoping property in the idealized case. Defining the complementary
residual $\delta_t^{\text{post}} = r^t_{\text{random}} + \gamma(1{-}d_t)V_{\text{old}}(\omega^{t+1}) - \tilde V(\tilde\omega^t)$,
one has $\delta_t^{\text{pre}} + \delta_t^{\text{post}} = $ the ordinary PPO TD residual
exactly; PDS-GAE traces only $\delta_t^{\text{pre}}$. With the \emph{approximate}
critics used in practice, the resulting bias/variance trade-off is an empirical
question, evaluated in Sec.~\ref{sec:results}.
\end{remark}

\begin{algorithm}[t]
\caption{PPO + Post-Decision State (PDS-GAE)}
\label{alg:pdsgae}
\begin{algorithmic}[1]
\State Initialize actor $\pi_\theta$, critics $V,\tilde V$, multipliers $\mu_D,\mu_P,\mu_B$
\For{rollout $=1,2,\dots$}
  \For{$t=0,\dots,n_{\text{steps}}-1$}
    \State Observe $\omega^t$; predict viewport; sample $\alpha^t \sim \pi_\theta(\cdot|\omega^t)$
    \State Execute $\alpha^t$, form $\tilde\omega^t$; observe $V^t_{\text{actu}}, h^{t+1}$
    \State Compute $r^t$, split into $r^t_{\text{known}}, r^t_{\text{random}}$
    \State Store $(\omega^t,\alpha^t,\tilde\omega^t,r^t_{\text{known}},r^t_{\text{random}},\omega^{t+1},d^t,V_{\text{old}}(\omega^t))$
  \EndFor
  \For{each stored step $t$} compute $\tilde y^t$, $\delta_t^{\text{pre}}$ \EndFor
  \State Backward per episode: compute $\hat A_t^{\text{PDS-GAE}}$, $\hat R_t^{\text{PDS-GAE}}$
  \State Update $\pi_\theta$ with $\hat A_t^{\text{PDS-GAE}}$; fit $V \to \hat R_t^{\text{PDS-GAE}}$, $\tilde V \to \tilde y^t$
  \State Update $\mu_D,\mu_P,\mu_B$ once per rollout (dual ascent)
\EndFor
\end{algorithmic}
\end{algorithm}

\subsection{Virtual Experience (proposed, not yet evaluated)}
\label{sec:ve}

The actual viewport $V^t_{\text{actu}}$ and next channel gain $h^{t+1}$ are conditionally
independent of a \emph{hypothetical} action given the current state -- in our
environment this holds by construction: $h^{t+1}$ is drawn i.i.d.\ from a fixed-parameter
Rician distribution regardless of the state or action, and $V^t_{\text{actu}}$ replays
from a pre-recorded head-movement trace. Consequently, the single realized pair
$(V^t_{\text{actu}}, h^{t+1})$ can be replayed, unchanged, against $M_v$ hypothetical
actions at the same timestep to generate $M_v$ additional $(\tilde\omega^{t,(m)}, \tilde
y^{t,(m)})$ training pairs for $\tilde V$ alone, using only the already-deterministic
physics functions (data-volume, transmission-rate, buffer, and coverage computations) --
no extra environment interaction.

\begin{algorithm}[t]
\caption{PDS + Virtual Experience}
\label{alg:ve}
\begin{algorithmic}[1]
\State Initialize $\tilde V$; choose VE batch size $B_{\text{VE}}$, update period $T$
\For{$t=0,1,2,\dots$}
  \State Observe $\omega^t$; select $\alpha^t=(\mbf{x}^t,P^t)$; form realized $\tilde\omega^t$
  \State Observe $V^t_{\text{actu}}, h^{t+1}$; compute $r^t, \omega^{t+1}$
  \If{$t \bmod T = 0$}
    \State Construct $\mathcal{B}_{\text{VE}}^t = \{\tilde\omega^{t,(1)},\dots,\tilde\omega^{t,(B_{\text{VE}})}\}$
    \For{each $\tilde\omega^{t,(b)} \in \mathcal{B}_{\text{VE}}^t$}
      \State Reuse $V^t_{\text{actu}}, h^{t+1}$; compute virtual $r^{t,(b)}, \omega^{t+1,(b)}$
      \State Update $\tilde V(\tilde\omega^{t,(b)})$
    \EndFor
  \Else
    \State Update only the realized $\tilde V(\tilde\omega^t)$
  \EndIf
\EndFor
\end{algorithmic}
\end{algorithm}

The scheme by which the $B_{\text{VE}}$ hypothetical actions are chosen is an open
design decision we have not yet fixed: sampling near $\pi_\theta(\cdot|\omega^t)$
concentrates the extra targets where the policy's own decision boundary lives, while
uniform sampling over the (combinatorially large, $\mathcal{O}(2^{50})$ effective free
tile choices at our tile budgets) action space would dilute the benefit. We report no
empirical results for VE in this paper; Sec.~\ref{sec:discussion} discusses why, given
our diagnostic findings, we do not expect it to address the feasibility oscillation
reported in Sec.~\ref{sec:results}.

% =====================================================================
\section{Experimental Setup}

The environment tiles a $360^{\circ}{\times}180^{\circ}$ panorama into an $8{\times}8$
grid ($45^{\circ}{\times}22.5^{\circ}$ tiles); the actual-viewport rectangle is
$120^{\circ}{\times}60^{\circ}$, forcing $9$--$16$ mandatory tiles per step. Episodes are
$36$ steps; one rollout is $48$ episodes ($1728$ environment steps), $\gamma{=}0.99$.
We evaluate two budgets, both with the power constraint \emph{genuinely removed}
(multiplier pinned at $0$ for the run's duration and excluded from the feasibility
check, not merely loosened): Budget A, $(\bar D,\bar B){=}(0.3,8.0)$; Budget B,
$(\bar D,\bar B){=}(0.1,8.0)$. Each budget/method pair is trained from multiple
independent seeds (Budget A: 2 seeds mature at matched rollout count at the time of
writing, additional seeds in progress; Budget B: 3 seeds). Held-out evaluation uses
$3$ fixed trace/seed combinations, deliberately \emph{not} refreshed per evaluation
call, so that all training seeds and both methods are compared against identical test
conditions. Following our own internal metrics-of-record convention, feasibility
(constraint satisfaction) is the primary comparison metric; quality among feasible
rollouts is secondary; raw Lagrangian reward is not used for cross-run comparison,
since it embeds the (method-specific) multiplier trajectory.

% =====================================================================
\section{Numerical Results}
\label{sec:results}

\subsection{Early convergence speed}

Table~\ref{tab:speed} reports, per seed, the first rollout at which a rolling-$20$-window
mean of the violation metric $\text{viol} = [J_D-\bar D]_+ + [J_B - \bar B]_+$ drops
below a given threshold and stays there. At the two loosest thresholds, PDS-GAE is
faster than the matched PPO seed in \emph{every} seed tested, in both budgets
(speedups of $2.1$--$3.5\times$). This advantage persists to near-exact satisfaction in
Budget A (every PDS seed faster than every PPO seed), but is contested in Budget B:
two of three PPO seeds reach near-zero violation faster than all three PDS-GAE seeds,
with the mean favoring PDS-GAE only because of one markedly slow PPO seed. We report
both the mean and the per-seed values for this reason -- the mean alone is not a
robust summary at $n{=}3$.

\begin{table}[t]
\centering
\caption{Rollouts to sustained constraint violation below threshold, per seed.}
\label{tab:speed}
\begin{tabular}{@{}llrrrr@{}}
\toprule
Budget & Method & viol$<5$ & viol$<1$ & viol$\approx 0$ \\
\midrule
A & PPO      & 549, 528           & 1166, 1152        & 1920, 2020 \\
A & PDS-GAE  & \textbf{173, 250}  & \textbf{541, 512}  & \textbf{1432, 1574} \\
B & PPO      & 576, 731, 499      & 849, 1597, 735     & 975, \emph{3269}, 866 \\
B & PDS-GAE  & \textbf{187, 232, 143} & 600, 465, 330  & 1281, 1515, 1477 \\
\bottomrule
\end{tabular}
\end{table}

\subsection{Mechanism: a measurably lower-variance advantage residual}

For every PDS-GAE rollout, we compute a zero-extra-cost counterfactual: the ordinary TD
residual on the identical rollout, using arrays already in scope (no extra network
forward passes). Pooled across all evaluated PDS-GAE arms ($\approx\!20{,}000$
rollout-rows), the PDS residual's standard deviation is lower than the ordinary
counterfactual's in $74.3\%$ of rollouts (mean ratio $0.92$, median $0.91$); the
effect is concentrated early in training (ratio $\approx\!0.65$--$0.85$ in the first
$200$ rollouts) and shrinks as training progresses (ratio $\approx\!0.93$--$0.98$ late).
This is directly consistent with Table~\ref{tab:speed}'s pattern: the speed advantage is
largest and most robust at loose/early convergence and weakens toward strict,
late-training satisfaction. We note, however, that the correlation between the PDS and
ordinary advantage estimates is weak ($0.11$--$0.28$) and their sign agreement is barely
above chance ($53$--$55\%$) -- the PDS residual is a lower-variance estimate of a
signal that is only partly the same as the ordinary one, not simply a cleaner version
of it.

\subsection{Final reliability and quality}

Table~\ref{tab:final} reports tail-window (last $200$ rollouts) statistics. PPO reaches
higher final feasibility than PDS-GAE in both budgets. Quality among feasible rollouts
(PSNR, and the discounted quality return $J_Q$) does not favor either method
consistently across budgets: PDS-GAE is slightly ahead in Budget B, PPO slightly ahead
in Budget A -- we report both rather than the single-budget snapshot that would
otherwise overstate a general quality advantage.

\begin{table}[t]
\centering
\caption{Tail-window (last 200 rollouts) comparison.}
\label{tab:final}
\begin{tabular}{@{}llrrr@{}}
\toprule
Budget & Method & Feasible \% & PSNR (dB)$^\dagger$ & $J_Q^\dagger$ \\
\midrule
A & PPO      & \textbf{100.0} & \textbf{51.92} & \textbf{25.58} \\
A & PDS-GAE  & 69.8  & 51.88 & 25.47 \\
B & PPO      & \textbf{94.8}  & 51.87 & 25.42 \\
B & PDS-GAE  & 59.7  & \textbf{51.94} & \textbf{25.58} \\
\bottomrule
\end{tabular}
\vspace{2pt}
\raggedright{\footnotesize $^\dagger$ Computed over feasible rollouts only.}
\end{table}

\subsection{A persistent, reproducible feasibility oscillation}

Beyond the tail-window snapshot, the full training trajectory (Fig.~\ref{fig:violation})
shows a qualitative difference between methods: PPO, once it first reaches feasibility,
settles into long, stable plateaus (in Budget A, a flat $100\%$ for over $2{,}000$
consecutive rollouts with no dips). PDS-GAE reaches high feasibility \emph{earlier} in
both budgets but never stabilizes the same way -- it oscillates repeatedly for the
remainder of training, independently reproduced in both budgets over $4{,}000+$
rollouts. This is, to our knowledge, a genuine and reproducible property of PDS-GAE in
this setting, not a transient one-off relapse.

\begin{figure}[t]
\centering
\includegraphics[width=\linewidth]{../FINAL_PART1_EXP/official_multiseed/A_dropP/01_violation.png}
\caption{Constraint violation and rolling feasible-\% over the full training
trajectory, Budget A (mean $\pm$ 1 std across seeds). PPO stabilizes after first
reaching feasibility; PDS-GAE does not.}
\label{fig:violation}
\end{figure}

% =====================================================================
\section{Discussion and Limitations}
\label{sec:discussion}

\textbf{What the oscillation is not.} We tested, and did not find support for, an
actor-update-aggressiveness explanation: at four independent relapse events, PPO's own
training diagnostics (approximate KL divergence, policy entropy, clip fraction) were
flat across the transition into relapse, not spiking. The PDS advantage magnitude
(std\_adv\_pds) also did not shrink during the periods of unresponsiveness we examined,
ruling out a vanishing-training-signal explanation.

\textbf{Partial evidence found.} (i) The policy's discretionary tile-count behavior
responded only slowly to a correctly, monotonically rising tile-budget multiplier near
the constraint boundary, in one case remaining flat for $600{+}$ rollouts despite
continued multiplier growth -- by contrast, the PPO baseline's tile count tracked its
own rising multiplier closely and promptly in the same cell. (ii) A genuine (not
sampling-noise, since evaluation seeds are fixed) gap between training-side and
held-out constraint costs on the stall dimension, whose breadth we cannot yet attribute
to broad generalization failure versus idiosyncrasy of the (necessarily small, $n{=}3$)
held-out trace set. (iii) In the budget with the larger feasibility gap, PDS-GAE's
policy entropy remained persistently higher late in training than most PPO seeds',
consistent with -- though not proof of -- incomplete convergence to a deterministic
policy.

\textbf{Why we do not evaluate VE in this paper.} VE (Sec.~\ref{sec:ve}) targets
sampling efficiency of the PDS critic's own training signal. Given (a) the measured
residual-variance advantage already shrinks on its own as real training data
accumulates, and (b) none of the mechanisms above implicate critic sample-starvation as
the oscillation's cause, we do not expect VE, even if implemented exactly as specified,
to resolve the reliability gap reported in Sec.~\ref{sec:results}. We regard it as a
theoretically well-motivated, cheaply implementable extension worth evaluating
separately, not a fix for the finding in this paper.

\textbf{Limitations.} Results use two budgets and $2$--$3$ mature seeds per
arm; additional seeds were in progress at submission time. % TODO: update seed counts / add
% new results here once the additional seeds mature -- see paper/NOTES.md
The held-out evaluation set is necessarily small (3 traces); broadening it is planned
future work, as is a targeted experiment raising the constraint-multiplier learning
rate to test the slow-responsiveness hypothesis above directly.

% =====================================================================
\section{Conclusion}

We adapted the post-decision-state critic mechanism to PPO for a constrained UAV
360\textdegree{} streaming CMDP, showed empirically and diagnostically that its
originally-proposed one-step advantage truncates credit assignment, and derived
PDS-GAE to address this. PDS-GAE reaches moderate constraint satisfaction
substantially faster than a matched PPO baseline in every seed tested, an effect we
connect directly to a measured reduction in advantage-estimator variance early in
training. This advantage does not yet extend to final reliability: PDS-GAE shows a
persistent, reproducible feasibility oscillation that we characterize but do not fully
resolve here. We proposed Virtual Experience as a theoretically motivated extension and
argued, on diagnostic grounds, why it is unlikely to be a fix for this specific gap,
leaving its empirical evaluation, along with a broader multiplier-learning-rate study
and a wider held-out trace set, to future work.

% =====================================================================
\begin{thebibliography}{9}

\bibitem{ref_mmsp25}
% TODO: replace with the real MMSP'25 baseline-paper citation (authors, title, venue, year).

\bibitem{ref_ehs}
% TODO: replace with the real EHS SAC+PDS paper citation.

\bibitem{felizardo2025pds}
G.~S.~Felizardo \emph{et al.}, ``A Reinforcement Learning Method for Environments with
Stochastic Variables: Post-Decision Proximal Policy Optimization with Dual Critic
Networks,'' \emph{arXiv:2504.05150}, 2025.

\bibitem{powell_adp}
W.~B.~Powell, \emph{Approximate Dynamic Programming: Solving the Curses of
Dimensionality}, 2nd ed. Wiley, 2011.

\bibitem{schulman2016gae}
J.~Schulman, P.~Moritz, S.~Levine, M.~Jordan, and P.~Abbeel, ``High-Dimensional
Continuous Control Using Generalized Advantage Estimation,'' \emph{arXiv:1506.02438},
2016.

\bibitem{schulman2017ppo}
J.~Schulman, F.~Wolski, P.~Dhariwal, A.~Radford, and O.~Klimov, ``Proximal Policy
Optimization Algorithms,'' \emph{arXiv:1707.06347}, 2017.

\bibitem{hca_harutyunyan}
A.~Harutyunyan \emph{et al.}, ``Hindsight Credit Assignment,'' \emph{NeurIPS}, 2019.
% TODO: verify exact venue/arXiv id before submission.

\bibitem{mesnard2021cca}
T.~Mesnard \emph{et al.}, ``Counterfactual Credit Assignment in Model-Free
Reinforcement Learning,'' \emph{arXiv:2011.09464}, 2021.

\end{thebibliography}

\end{document}
